 \documentclass{article}
\usepackage{amsmath}
\usepackage{graphicx}
\graphicspath{ {images/} }
\usepackage[spanish]{babel}
\usepackage{bbm}
\usepackage{amsthm}
\usepackage{authblk}
\usepackage{hyperref}
\usepackage[utf8]{inputenc}
\usepackage{mathrsfs}
\usepackage{amsfonts}
\usepackage{amssymb}
\usepackage{enumerate}
\usepackage[left=3cm,right=2cm,top=2cm,bottom=2cm]{geometry}
\usepackage{epsfig}
 \renewcommand{\baselinestretch}{0.93}
 \thispagestyle{empty}
 \pagestyle{empty}
\setlength{\textheight}{53pc} \setlength{\textwidth} {36pc}
\setlength{\topmargin}{-4mm}
\usepackage{multicol}
\newcommand*{\email}[1]{%
    \normalsize\href{mailto:#1}{#1}\par
    }
\setlength{\columnsep}{1cm}
\newcommand{\N}{\mathbb{N}}
\newcommand{\calM}{\cal{M}}
\newcommand{\calP}{\cal{P}}
\newcommand{\F}{\mathbb{F}}
\newcommand{\R}{\mathbb{R}}
\newcommand{\Z}{\mathbb{Z}}
\newcommand{\Q}{\mathbb{Q}}
\newcommand{\C}{\mathbb{C}}
\newcommand{\bbH}{\mathbb{H}}


\theoremstyle{definition}
\newtheorem{ejer}{Ejercicio}
\author{Marcos Morales}
\affil{Universidad Finis Terrae}
\title{Primera Semana\\}



	


\begin{document}


\begin{picture}(400, 75)
   \put(-15,15){\includegraphics[width=5cm]{UFT.png}}
\end{picture}



\begin{center}
\textbf{\large{Laboratorio Evaluado 2}} \\
\end{center}

\setlength{\parindent}{0cm}



\begin{center}
\textbf{Cálculo Vectorial}
\end{center}


\begin{center}
\textbf{\large{Análisis del Modelo}} \\
\end{center}



\hspace{100cm} \\ 


Considere que los dos modelos Cobb-Douglas de dos empresas $A$ y  $B$ están dados por

\begin{align*}
    Y_A &= 10K^{0.3}L^{0.7} \\
    Y_B &= 7K^{0.83}L^{0.17}
\end{align*}

El objetivo de este laboratorio es que se realize un análisis exaustivo del modelo Cobb-Douglas, comparando las producciones de estas dos empresas. \\

\begin{center}
    \textbf{\Large{Actividades}}
\end{center}

\begin{itemize}
    \item Usando Matlab u Ocatve, realice un gráfico de ambos modelos.
    \item Los modelos presentan máximos relativos?, mínimos relativos?, puntos silla?. Argumente usando el gradiente de la función. 
    \item Compare los modelos en distintos puntos del plano, incluya al menos 10 puntos.
    \item Considere que el factor capital y la dotación del trabajo están relacionados por \\$K+L=C$, donde $C$ es una constante. Determine  los puntos extremos con esa condición
    \item ¿Existe algún valor de $C$ para el cual el máximo relativo del punto anterior de $Y_A$ sea mayor al de $Y_B$ ?. En caso de ser cierto determine los valores posibles.
    \item Explique la razón matemática por la cuál ocurre esto 
    \item Considere ahora $Y_D = 7K^{0.5}L^{0.5}$ e $Y_E = 7K^{0.5}L^{0.5}$, y realice la comparación de los máximos relativos de $Y_A$, $Y_B$, $Y_D$ e $Y_E$
\end{itemize}









\newpage


\begin{center}
\textbf{Aspectos Formales}
\end{center}

- Debe presentar un informe de a lo más 8 páginas con introducción, desarrollo y conclusión. \\

- Usted mismo debe explicar qué significa cada variable en las ecuaciones diferenciales, explicando cual es el significado físico de las mismas. \\

- Los grupos deben ser de a lo más 5 integrantes. \\

- Debe resolver las ecuaciones diferenciales de forma analítica. \\

- Debe resolver las ecuaciones diferenciales de forma computacional utilizando Matlab u Octave. \\

- Debe incluir los gráficos de las soluciones obtenidas.\\

- Los códigos Matlab usados deben ser entregados en un documento aparte. \\

- En la sección de conclusión debe incuir una breve discusión en la cuál expliquen las reacciones de cada integrante del grupo mientras desarrollaban el laboratorio. \\












\end{document}